% ======================================================================
% SECTION 3: AUTONOMOUS TRIAL DESIGN AND STATISTICAL FRAMEWORK
% File: sections/trial_design.tex
% ======================================================================

Autonomous trial design replaces the traditional interplay between clinical scientists,
biostatisticians, and regulatory strategists with integrated agents that generate protocol
documents, select design patterns, define endpoints and estimands, and produce statistical
analysis plans. The clinical accountability agent provides medical logic and dose optimization,
while the biostatistics and data agent handles statistical planning and interim analysis
design. Together, these agents produce ICH M11-compliant structured protocols \cite{ich-m11}
with embedded adaptive decision rules and Bayesian analysis specifications aligned with the
FDA January 2026 draft guidance on Bayesian methods \cite{fda-bayesian-guidance}.

\subsection{Clinical Accountability Agent}

The clinical accountability agent (clinical\_accountability\_agent) replaces the study
responsible physician and clinical lead. It automates protocol medical logic, eligibility
criteria generation, dose and schedule optimization, and medical monitoring rule definition.

Eligibility criteria generation operates through a rule engine that translates indication
characteristics, biomarker requirements, prior therapy constraints, and safety
considerations into structured inclusion and exclusion criteria. The engine references the
adapted ICH E6(R3) protocol requirements \cite{ich-adapt}, which specify that eligibility
criteria must balance scientific rigor with practical enrollment feasibility. For
biomarker-driven trials, the agent integrates biomarker prevalence data and testing
access estimates to ensure that eligibility criteria do not create enrollment bottlenecks.

Dose and schedule optimization follows FDA Project Optimus requirements
\cite{fda-project-optimus}, which mandate characterization of the dose-exposure-response
relationship before advancing to registrational trials. The agent designs dose-finding
schemas using model-based approaches (including modified continual reassessment method,
Bayesian optimal interval designs such as BOIN, and modified toxicity probability
interval designs) that characterize both safety and efficacy across the dose range. For
combination therapies and newer modalities such as antibody-drug conjugates and bispecific
antibodies, the agent incorporates modality-specific dose optimization considerations
including target-mediated drug disposition, immunogenicity, and schedule-dependent
pharmacodynamics.

Medical monitoring rules define the ongoing safety oversight framework for the trial.
The agent generates rules for routine safety review frequency, dose modification
criteria, treatment discontinuation triggers, and safety stopping boundaries. These
rules are encoded as machine-readable policies that integrate with the safety agent
(\S\ref{sec:safety}) for real-time monitoring execution.

\subsection{Autonomous Protocol Design}

The autonomous protocol design subsystem generates structured protocol documents following
the ICH M11 clinical electronic structured harmonised protocol standard \cite{ich-m11}.
The system selects from six validated design patterns based on indication characteristics,
biomarker strategy, regulatory pathway, and operational feasibility constraints.
Table~\ref{tab:design-patterns} compares these six design patterns.

\begin{longtable}{L{2.0cm} L{2.2cm} L{2.5cm} L{2.2cm} L{2.8cm}}
\caption{Trial Design Pattern Comparison: Six Validated Design Approaches}
\label{tab:design-patterns} \\
\toprule
\textbf{Design Pattern} & \textbf{Best For} & \textbf{Key Advantage} & \textbf{Regulatory Basis} & \textbf{Agent Configuration} \\
\midrule
\endfirsthead
\toprule
\textbf{Design Pattern} & \textbf{Best For} & \textbf{Key Advantage} & \textbf{Regulatory Basis} & \textbf{Agent Configuration} \\
\midrule
\endhead
\midrule
\multicolumn{5}{r}{\textit{Continued on next page}} \\
\endfoot
\bottomrule
\endlastfoot
Conventional RCT & Well-characterized indications with large populations & Highest interpretability, gold standard evidence & FDA Guidance on Clinical Trials \cite{fda-clinical-trials-guidance} & Standard randomization, fixed sample size, parallel arms \\
Seamless Phase I/II or II/III & Programs seeking cycle-time reduction & Eliminates inter-phase gaps, continuous learning & ICH E20 Adaptive Design \cite{ich-e20} & Adaptive boundaries, interim decision rules, information fraction triggers \\
Basket Master Protocol & Rare biomarker-defined populations across tumor types & Cross-tumor biomarker evaluation, shared infrastructure & FDA Master Protocol Guidance \cite{fda-adaptive-designs} & Biomarker-first enrollment, tumor-agnostic analysis, borrowing algorithms \\
Umbrella Master Protocol & Multiple biomarker-defined subgroups within a single tumor & Efficient within-tumor screening, shared control arm & FDA Master Protocol Guidance \cite{fda-adaptive-designs} & Screening platform, sub-study randomization, shared infrastructure \\
Platform Trial & Long-running evaluation programs with arms added over time & Shared control, perpetual infrastructure & ICH E20 Adaptive Design \cite{ich-e20} & Non-concurrent control adjustment, arm addition and graduation rules \\
Single-arm with External Control & Rare diseases or settings where RCT is infeasible & Speed and feasibility for unmet need & FDA RWD/RWE Framework \cite{fda-rwd-registries} & RWE source selection, propensity matching, sensitivity analysis \\
\end{longtable}

The design selection algorithm evaluates each candidate design against five criteria:
patient population size and prevalence (favoring master protocols for rare biomarkers),
regulatory pathway alignment (favoring conventional RCT for standard approval and
single-arm for accelerated approval), operational complexity tolerance, competitive
timing pressure, and available infrastructure. The algorithm produces a ranked design
recommendation with justification for each criterion.

Adaptive design governance follows ICH E20 principles \cite{ich-e20} and FDA guidance
on adaptive designs \cite{fda-adaptive-designs}. For trials incorporating interim
analyses, sample size re-estimation, or treatment arm modifications, the agent
generates governance procedures that specify access controls for accumulating data,
documentation requirements for each adaptive decision, simulation artifacts
demonstrating operating characteristics, and independence safeguards between the
analysis team and the study conduct team. These governance procedures ensure that
adaptive modifications maintain trial integrity while enabling efficient use of
accumulating evidence.

\subsection{Biomarker and Companion Diagnostics Strategy}

The biomarker strategy subsystem manages the biomarker development lifecycle from
analytical validation through clinical utility demonstration and companion diagnostic
co-development. This subsystem addresses a critical operational challenge: biomarker-driven
oncology trials increasingly depend on companion diagnostics that must be developed in
parallel with the therapeutic, creating regulatory and operational interdependencies.

The biomarker evidence hierarchy guides development decisions through three stages.
Analytical validity establishes that the assay reliably measures the biomarker of
interest with acceptable precision, reproducibility, and robustness. Clinical validity
demonstrates that the biomarker measurement correlates with a clinical outcome or
biological process relevant to the therapeutic hypothesis. Clinical utility proves that
use of the biomarker in clinical decision-making improves patient outcomes compared with
not using the biomarker.

Companion diagnostic co-development planning addresses the parallel regulatory pathways
for the therapeutic and the diagnostic. In the United States, the agent coordinates IND
and IDE submissions when the companion diagnostic requires investigational device
exemption. In the European Union, the agent tracks EU In-Vitro Diagnostic Regulation
(IVDR) requirements that affect the availability and regulatory status of the diagnostic.
The agent generates integrated development timelines that synchronize therapeutic and
diagnostic milestones.

Circulating tumor DNA integration as a monitoring biomarker represents a specific
application of the biomarker strategy. The agent tracks FDA guidance on ctDNA as a
biomarker for treatment response monitoring, molecular residual disease detection, and
early relapse identification. For Physical AI trials incorporating robotic surgical
procedures, ctDNA monitoring provides a molecular complement to imaging and physical
examination for post-surgical surveillance.

\subsection{Biostatistics and Data Agent}

The biostatistics and data agent (data\_biostats\_agent) replaces the biostatistics team for
statistical planning, analysis execution, and regulatory submission dataset preparation.
The agent generates statistical analysis plans with Bayesian adaptive methods, designs
interim analyses with futility and efficacy boundaries, implements adaptive randomization
logic, and ensures CDISC SDTM and ADaM compliance \cite{cdisc-sdtm, cdisc-adam}.

The ICH E9(R1) estimand framework \cite{ich-e9r1} provides the foundational structure for
defining treatment effects in a manner that is aligned across clinical, statistical, and
regulatory perspectives. Table~\ref{tab:estimand} presents the five estimand components
and their implementation in the autonomous system.

\begin{longtable}{L{2.2cm} L{3.0cm} L{3.0cm} L{3.5cm}}
\caption{Estimand Framework Components per ICH E9(R1)}
\label{tab:estimand} \\
\toprule
\textbf{Component} & \textbf{Definition} & \textbf{Agent Implementation} & \textbf{Oncology Example} \\
\midrule
\endfirsthead
\toprule
\textbf{Component} & \textbf{Definition} & \textbf{Agent Implementation} & \textbf{Oncology Example} \\
\midrule
\endhead
\midrule
\multicolumn{4}{r}{\textit{Continued on next page}} \\
\endfoot
\bottomrule
\endlastfoot
Treatment & The treatment condition of interest and the alternative & Protocol design module specifies treatment arms, doses, schedules, and comparators & Pembrolizumab 200mg Q3W versus investigator choice chemotherapy \\
Population & The patients targeted by the clinical question & Eligibility engine defines inclusion and exclusion criteria with biomarker requirements & Adults with PD-L1 TPS $\geq$ 50\% advanced NSCLC, no prior systemic therapy \\
Variable (endpoint) & The outcome variable that is measured & Endpoint module selects primary, secondary, and exploratory endpoints & Overall survival, measured as time from randomization to death from any cause \\
Intercurrent events & Events occurring after treatment initiation that affect interpretation & Event classification engine categorizes treatment discontinuation, rescue therapy, and death & Treatment switch handled via treatment policy strategy; study drug discontinuation via hypothetical strategy \\
Population-level summary & The method used to summarize the treatment effect & Analysis engine implements hazard ratio, restricted mean survival time, or response rate & Hazard ratio from stratified log-rank test with pre-specified strata \\
\end{longtable}

Bayesian adaptive methods follow the FDA January 2026 draft guidance
\cite{fda-bayesian-guidance}, which normalizes the use of Bayesian approaches in
clinical trials. The agent implements explicit prior specification with justification
(informative priors from historical data, weakly informative priors for novel settings,
or non-informative priors when historical data is limited), robustness and sensitivity
analysis across prior specifications, simulation-based operating characteristics
validation, and clear success criteria defined in terms of posterior probabilities.
For interim analyses, the agent computes posterior predictive probabilities of trial
success to support futility stopping decisions, and uses Bayesian decision rules for
adaptive randomization that allocate patients preferentially to better-performing arms
while maintaining adequate data for comparative inference.
